RailFlux Database Architecture - Technical Report
Three-Schema Separation Architecture
Schema Segregation: PostgreSQL implementation utilizes three distinct schemas providing clear separation of concerns and operational boundaries. Railway_config contains immutable configuration data including signal_types, signal_aspects, and point_positions defining operational parameters. Railway_control manages live operational data encompassing track_circuits, track_segments, signals, point_machines, interlocking_rules, and system_state tables. Railway_audit provides comprehensive audit trail functionality through event_log, system_events, and audit_summary tables ensuring regulatory compliance and forensic analysis capabilities.
Cross-Schema Referential Integrity: Foreign key constraints span schema boundaries establishing data consistency between configuration, operational, and audit domains. Signal records reference signal_types across schema boundaries while maintaining transactional consistency through PostgreSQL's ACID compliance.
Core Operational Tables Structure
Track Circuit Hardware Layer: track_circuits table serves as physical hardware abstraction with circuit_id as primary identifier, is_occupied boolean state, occupied_by train identification, last_occupied_at timestamps, failure states, and protecting_signals array establishing signal-to-circuit protection relationships. GIN indexes on protecting_signals array enable efficient reverse lookups for safety validation.
Track Segment Virtual Layer: track_segments table provides logical abstraction over physical circuits with segment_id business identifiers, visual coordinate pairs (start_row, start_col, end_row, end_col), circuit_id foreign key linking to hardware layer, is_assigned logical occupancy state, protecting_signals array, length_meters physical measurements, and max_speed_kmh operational constraints. This dual-layer approach separates physical hardware detection from logical operational sections.
Signal Control Infrastructure: signals table encompasses comprehensive signal management with signal_id business keys, signal_name display identifiers, signal_type_id references to configuration lookup, location coordinates for visual rendering, direction enumeration (UP/DOWN), current_aspect_id references to configuration aspects, aspect_count defining signal capabilities, possible_aspects array defining permitted indications, protected_track_circuits array establishing protection relationships, interlocked_with integer arrays for signal dependencies, activation states, and audit timestamp fields.
Point Machine Mechanical Control: point_machines table manages switching infrastructure with machine_id identifiers, position_id references to configuration lookup, junction coordinates, current_position_id state tracking, operating_status enumeration (CONNECTED, IN_TRANSITION, FAILED, LOCKED, MAINTENANCE), lock states, transition timing, operation counts, safety_interlocks arrays, protected_signals arrays, paired_entity references for coordinated machines, and route assignment extensions supporting automatic normalization and locking capabilities.
Route Assignment Extension Tables
Track Circuit Graph Structure: track_circuit_edges table implements A* pathfinding graph with from_circuit_id and to_circuit_id establishing directed connections, side enumeration (LEFT/RIGHT) for bidirectional tracks, condition_point_machine_id and condition_position implementing conditional edge traversal based on point machine states, weight values for pathfinding optimization, and activation flags enabling dynamic graph topology management.
Route State Management: route_assignments table provides comprehensive route lifecycle tracking with UUID primary keys, source_signal_id and dest_signal_id establishing route endpoints, direction constraints, assigned_circuits arrays containing complete route paths, overlap_circuits arrays for safety margins, state enumeration tracking route progression through REQUESTED, VALIDATING, RESERVED, ACTIVE, PARTIALLY_RELEASED, RELEASED, FAILED, EMERGENCY_RELEASED, DEGRADED states, timing fields for lifecycle management, locked_point_machines arrays, priority values, operator identification, failure reasons, and performance_metrics JSONB containing execution statistics.
Resource Conflict Management: resource_locks table implements comprehensive locking mechanism with UUID identifiers, resource_type enumeration (TRACK_CIRCUIT, POINT_MACHINE, SIGNAL), resource_id string references, route_id foreign keys, lock_type enumeration (EXCLUSIVE, SHARED, OVERLAP, EMERGENCY), acquisition timestamps, operator tracking, expiration management, activation flags, release tracking, and release reasons enabling complete resource conflict detection and resolution.
Signal Safety Overlaps: signal_overlap_definitions table manages safety braking distances with signal_id references, overlap_circuits arrays defining protected regions, overlap_distance_m measurements, release_conditions arrays, overlap_type enumeration (FIXED, VARIABLE, FLANK_PROTECTION), hold_seconds timing parameters, and activation management ensuring train protection beyond primary route boundaries.
Real-Time Notification System
PostgreSQL LISTEN/NOTIFY Implementation: Database triggers automatically generate notifications using pg_notify function calls with railway_changes channel for operational updates and route_changes channel for route-specific events. Trigger functions construct JSON payload objects containing table names, operation types, entity identifiers, and timestamps enabling precise change tracking.
Application-Level Integration: DatabaseManager implements QSqlDriver notification subscription with signal-slot connections translating database notifications into Qt signals. Real-time updates propagate through QML property bindings enabling immediate UI updates without polling overhead. Notification processing includes payload parsing, entity identification, and appropriate signal emission for UI component updates.
Performance Optimization: Dual-mode operation combines real-time notifications for immediate updates with configurable polling (400-500ms production intervals) providing fallback mechanism ensuring data consistency. Performance monitoring tracks notification latency, processing times, and queue depths maintaining sub-50ms response targets for safety-critical operations.
Advanced Database Functions
Atomic Operation Functions: PostgreSQL stored procedures implement complex atomic operations including update_signal_aspect with audit trail generation, update_point_position with paired machine coordination, simulate_track_occupancy for testing scenarios, acquire_resource_lock with conflict detection, release_route_resources with batch operations, and validate_route_request implementing comprehensive safety validation.
Trigger Function Ecosystem: Comprehensive trigger functions manage automatic operations including update_timestamp for generic modification tracking, update_signal_change_time for aspect-specific timing, set_event_date for audit partitioning, notify_railway_changes for real-time notifications, log_changes for comprehensive audit trails, and validate_interlocking_constraints for safety rule enforcement.
Conditional Logic Implementation: Advanced stored procedures implement railway-specific business logic including point machine pairing validation, track circuit aggregation for segment states, signal aspect progression validation, resource conflict resolution algorithms, and route state machine transitions with safety interlocks.
Safety Mechanisms and Compliance
Audit Trail Completeness: event_log table captures comprehensive operational history with event_timestamp precision, event_type categorization, entity_type and entity_id references, old_values and new_values JSONB objects enabling complete state reconstruction, operator identification, operation_source tracking, safety_critical boolean flags, authorization_level tracking, reason_code categorization, replay_data for forensic analysis, and sequence_number ordering ensuring complete audit trail coverage.
Safety-Critical Operation Validation: Database constraints implement railway safety rules including track circuit occupancy validation, signal aspect progression restrictions, point machine position consistency checking, resource lock conflict prevention, route state transition validation, timing constraint enforcement, and operator authorization verification ensuring operational safety through database-level enforcement.
Performance Monitoring Integration: Database functions implement performance tracking with response time measurement, operation success/failure rates, resource utilization monitoring, safety violation detection, compliance scoring, and automated alerting when thresholds exceed acceptable limits ensuring continuous safety and performance validation.
Index Strategy and Performance
Primary Performance Indexes: Comprehensive indexing strategy includes unique constraints on business keys, foreign key indexes for referential integrity, partial indexes on active records (WHERE is_active = TRUE), occupancy-specific indexes (WHERE is_occupied = TRUE), state-specific indexes for route management, and temporal indexes for audit queries enabling efficient query processing across all operational scenarios.
Advanced Index Types: GIN indexes on array columns (possible_aspects, protecting_signals, assigned_circuits) enable efficient array containment queries essential for safety validation. JSONB GIN indexes on audit fields enable flexible forensic analysis. Composite indexes on frequently joined columns optimize complex queries. Partial indexes reduce storage overhead while maintaining query performance.
Query Optimization Features: Database views provide optimized query interfaces including v_signals_complete combining signal data with configuration lookups, v_track_segments_active filtering operational segments, v_resource_utilization providing real-time resource usage statistics, and v_recent_events enabling efficient audit queries with automatic time-based filtering.
Connection Management and Reliability
Dual-Connection Architecture: DatabaseManager implements both system PostgreSQL connections (standard installation, port 5432) and portable PostgreSQL connections (portable installation, port 5433) with automatic failover providing deployment flexibility and reliability. Connection validation includes database accessibility testing, schema verification, function availability checking, and performance baseline establishment.
Error Recovery Mechanisms: Comprehensive error handling includes connection retry logic with exponential backoff, transaction rollback capabilities, deadlock detection and resolution, constraint violation handling, network interruption recovery, and graceful degradation modes ensuring continuous operation under adverse conditions.
Resource Management: Connection pooling prevents resource exhaustion, prepared statement caching improves performance, transaction scope management ensures data consistency, memory management prevents leaks during 24/7 operation, and connection lifecycle management handles application startup, operation, and shutdown scenarios.
Security and Access Control
Role-Based Access Control: Three-tier security model includes railway_operator role with full operational access, railway_observer role with read-only access, and railway_auditor role with audit trail access. Schema-level permissions ensure proper access boundaries. Function-level security prevents unauthorized operations. Row-level security potential for multi-tenant scenarios.
Audit and Compliance Features: Complete operation logging with operator identification, IP address tracking, session management, authorization level verification, reason code requirements for critical operations, safety-critical event flagging, regulatory compliance reporting, and forensic analysis capabilities meeting railway industry safety standards.
Data Integrity Protection: Foreign key constraints prevent orphaned records, check constraints enforce business rules, trigger-based validation implements complex safety rules, transaction isolation prevents concurrent modification issues, and backup/recovery procedures ensure data protection meeting railway operational requirements.